\renewcommand{\thechapterdisplay}{Vigésimo segundo dia}
\renewcommand{\thechapterrunner}{Terceiro modo: a Santa Comunhão}
\chapter[Vigésimo segundo dia --- Terceiro modo de aliviar as almas do Purgatório: a Santa Comunhão]{Terceiro modo de aliviar as almas do Purgatório: a Santa Comunhão}
\section*{Comunhão sacramental}

Quando temos a felicidade de comungar, nós nos unimos a Jesus Cristo de uma maneira tão íntima, que cada um de nós poderia exclamar como o grande apóstolo: \textit{Já não sou eu que vivo, mas é Cristo que vive em mim}. Então nossa carne se torna Sua própria carne, Seu coração se faz palpitar no nosso, Seu Sangue corre em nossas veias, Sua divindade reside em nós: somos outros Cristos! Ó, nestes felicíssimos momentos, invejados pelos anjos, como nos é fácil falar com Deus sem o ruido das palavras, e de dizer-lhe cheios de confiança: \textit{Ó Deus, protetor dos aflitos, olhai para mim e vereis a face de vosso Cristo. Respice in faciem Christi tui. Não sou mais eu que falo e que peço, é Jesus, vosso próprio filho, que fala e que pede por mim. É ele que roga por mim a libertação de meu pai, de minha mãe, a libertação dessas pobres almas abandonadas. Tenho certeza, ó Pai de misericórdia, que não rejeitareis estas justas súplicas, pois o rosto, as orações, as lágrimas, o Sangue de Jesus Cristo têm uma voz toda poderosa para apaziguar vossa justiça e obter o perdão!}

Comunguemos, portanto, alma cristã, comunguemos frequentemente por essas almas tão amadas que não têm mais a felicidade de participar do banquete Eucarístico. Ó, com que fervor, com que santa impaciência, elas esperam que nós espalhemos sobre elas o orvalho refrescante e libertador do Sangue de Jesus Cristo. Em breve --- como esse pensamento é consolador --- muito em breve, a comunhão eterna começará para eles, e eles contemplarão no Céu o Salvador eucarístico, que é o Deus das recompensas. 


\section*{Comunhão espiritual}

Se não podeis, alma cristã, fazer frequentemente a comunhão sacramental, ou seja, receber realmente a Jesus Eucarístico, fazei ao menos a comunhão espiritual. Ela consiste de um desejo ardente de se unir ao divino Salvador e de receber seu espírito e suas graças. É uma pratica muito proveitosa tanto para os viventes quanto para os mortos. Santo Afonso de Ligório chega ao ponto de dizer que podemos obter mais frutos da comunhão espiritual feita com fervor que da comunhão sacramental feita com tibieza. Também há a vantagem de que podemos fazê-la todos os dias, em qualquer momento do dia e da noite, e em qualquer lugar, seja profano ou sagrado. Não é um meio simples, fácil e poderoso de aliviarmos nossos queridos irmãos falecidos?

Fazei, portanto, alma cristã, a comunhão espiritual a cada visita ao Santíssimo Sacramento. Eis a fórmula que podeis empregar: \textit{Meu Jesus! Eu creio que estais presente aqui. Eu Vos amo, eu Vos desejo, eu me uno a Vós em espírito, de coração, esperando que eu possa vos receber realmente. Abençoai-me! Abençoai também as almas do Purgatório que sofrem tanto. Sim, Senhor, chamai Vossos filhos e nosso irmãos ao repouso eterno. Que a luz que não se apaga brilhe sobre eles! Que descansem em paz! Requiem æternam dona eis, Domine!}


\section*{Exemplo }

O Venerável Luis de Bois, célebre mestre da vida espiritual e homem de notável sabedoria, conta que um devoto servo de Deus, que ele conhecia e amava, foi visitado por uma alma do Purgatório, e que esta o fez contemplar todos os tormentos que padecia. Ela foi punida por ter recebido a divina Eucaristia com uma preparação insuficiente e muita tibieza. Por isso, tinha sido condenada ao suplício de um fogo devorador que a consumia. Ela disse: \textit{Eu vos imploro, vós que fostes meu amigo íntimo e fiel, e que ainda o é; eu vos imploro que comungueis uma vez em meu nome, com todo amor e devoção de que sois capaz. Tenho confiança de que essa comunhão fervorosa será suficiente para me libertar do Purgatório, e que por meio dela serão reparadas minhas culpáveis friezas.}

O homem se apressou em assistir à missa e em comungar piedosamente pelo repouso da alma conhecida. Após a ação de graças, a alma apareceu-lhe novamente, cercada de um brilho incomparável, feliz e cheia de gratidão: \textit{Sejais bendito, ó meu melhor amigo! Vossa comunhão me libertou e eu verei meu adorável Mestre face a face!} 

Convém agora relembrar o conselho de São Boaventura: \textit{Que a caridade vos leve à comungar frequente e piedosamente, pois não há nada mais eficaz para o repouso eterno dos defuntos.}


\vspace{\baselineskip}\noindent \textbf{Oremos.} Vós retendes, ó meu Deus, as almas de meus próximos em Vossa justiça. Mas vos quereis que, ao comer o Pão dos Anjos, eu possa abrir-lhes as portas do paraíso. Sejais bendito, Pai misericordioso, e recebei a promessa que Vos faço hoje de comungar frequentemente em favor dessas santas almas do Purgatório, afim de que não vejais mais em mim senão o Vosso Filho. Que minha voz, transformada na Sua, Vos alcance e me obtenha mais seguramente a graça que Vos peço. Ó Jesus, sede-lhes propício! Que descansem em paz!



\begin{center}
\textit{Requiescant in pace!}
\end{center}

Rezai agora uma dezena do terço e as orações no final deste livro.
